\section{Introduction}
\label{sec:introduction}

Transformer language models compute by applying a stack of residual blocks, and
every token pays for every block. Depth is therefore a direct axis of inference
cost: it multiplies latency, memory traffic, and energy, and unlike hidden width
it cannot be amortized away by batching. As models grow deeper to grow more
capable, the cost of simply traversing the stack has become one of the dominant
terms in the inference bill \citep{zhen2025taming}.

Layer skipping is an attractive response. Removing an entire block eliminates real
computation, its attention and MLP matmuls both disappear, without changing hidden
width, attention pattern or KV-cache layout, and a substantial literature establishes that
Transformers tolerate it surprisingly well \citep{men2024shortgpt, gromov2024unreasonable,
fan2020reducing}. Almost all of that work asks \emph{which} blocks to remove, and answers
with progressively better importance scores. Far less attention has gone to the operator
that does the removing. When a block is skipped, its residual update is replaced by zero.
That choice is never argued for; it is simply what ``removing a layer'' has come to mean.
It is computationally convenient, and representationally wasteful.

A separate line of work tells us the residual stream is not a sequence of unrelated states but a
trajectory: residual networks behave like ensembles of shallow paths \citep{veit2016residual},
residual connections encourage iterative refinement of a single representation
\citep{jastrzebski2018residual}, and a hidden state at one depth can often be mapped to a nearby
depth linearly \citep{din2023jump, belrose2023eliciting}. If depth is a trajectory, a block's
update is the next step along a path the model has already walked, and steps along a smooth path
should be predictable from the steps before them.

\begin{center}
\emph{How much of a skipped Transformer block's residual update is predictable from the preceding depth trajectory, \\ and can it be recovered at negligible
cost?}
\end{center}

We take this question literally. Writing the pre-norm update at layer $\ell$ as
$\dd_\ell = F_\ell(\hh_\ell)$, so that $\hh_{\ell+1} = \hh_\ell + \dd_\ell$, plain
skipping is precisely the estimator $\widehat{\dd}_\ell = 0$, a zero-order predictor that
ignores every update the model has already computed. We instead treat $\dd_{\ell-1},
\dd_{\ell-2}, \dots$ as a short series indexed by depth and fit an autoregressive model
that forecasts $\dd_\ell$ from it, using information already sitting in the residual
stream at inference time and adding 0.64\% to prefill latency (7B, sequence length
512).

Much of the missing update turns out to be recoverable. We introduce \textbf{\mname
(Depth-wise AutoRegressive update prediction)}, which replaces the zero-update assumption
with a closed-form per-channel forecast fitted on unlabeled text with no gradient
training, and we report three results.

Four things here are new. (1) \textbf{Per-channel depth-autoregressive recovery}: we forecast a
skipped block's residual update from the updates preceding it, per channel and in closed form,
and show that the obvious scalar version of the idea fails, and is at some layers worse than
dropping the update outright. (2) \textbf{A cross-token extension}, which also reads the previous
token's update from the preceding layer, and which repairs the one layer no amount of
regularisation could. (3) \textbf{Damage proportionality}: the benefit of repair scales with the
damage skipping inflicts on the likelihood axis, and significant downstream gains appear only
where that damage is severe,
an ordering we then bound with the first out-of-family test of it. (4) \textbf{A measured noise
floor} for multiple-choice evaluation: re-scoring an identical configuration under a different
batch shape moves the score, by as much as several of the effects reported in this area.

Concretely, \mname recovers 4.6\% to 23.3\% of the language-modelling damage skipping causes at
light budgets across Qwen2.5-0.5B, 1.5B and 7B, at 0.64\% added prefill latency. Where skipping
is destructive, that recovery reaches behaviour: at 28.6\% of blocks skipped on Qwen2.5-7B,
under the deployable rule, it answers 89 more questions correctly out of 900 than plain skipping
($z=4.31$), one of four results surviving Bonferroni correction over all 63 comparisons.
